% ============================================================================
% LANGENTWURF LE-002: Geräte-Inventur
% Profil: S (Senioren 65+) | Tier: 2 (Standard)
% ============================================================================

\documentclass[11pt,a4paper]{article}
\usepackage[utf8]{inputenc}
\usepackage[T1]{fontenc}
\usepackage[ngerman]{babel}
\usepackage{geometry}
\usepackage{fancyhdr}
\usepackage{lastpage}
\usepackage{xcolor}
\usepackage{tcolorbox}
\usepackage{enumitem}
\usepackage{tabularx}
\usepackage{longtable}
\usepackage{booktabs}
\usepackage{hyperref}
\usepackage{amssymb}

\geometry{left=2.5cm, right=2.5cm, top=2.5cm, bottom=2.5cm}
\definecolor{fach}{HTML}{b35c00}
\definecolor{gray}{HTML}{666666}
\definecolor{lightgray}{HTML}{f5f5f5}
\definecolor{successgreen}{HTML}{2a7a4b}
\definecolor{warningorange}{HTML}{b35c00}

\tcbuselibrary{skins,breakable}
\newtcolorbox{infobox}[1][]{ colback=lightgray, colframe=fach, fonttitle=\bfseries, title=#1, breakable }
\newtcolorbox{scorebox}{ colback=white, colframe=fach, fonttitle=\bfseries, title=Didaktik-Score, breakable }

\pagestyle{fancy}
\fancyhf{}
\fancyhead[L]{\small\textcolor{gray}{Digitale Grundlagen -- 002 Geräte-Inventur}}
\fancyhead[R]{\small\textcolor{gray}{LE Tier 2 / Profil S}}
\fancyfoot[C]{\small\thepage{} / \pageref{LastPage}}
\renewcommand{\headrulewidth}{0.4pt}

\begin{document}

% Deckblatt
\begin{titlepage}
    \centering
    \vspace*{2cm}
    {\Large\textcolor{gray}{Digitale Grundlagen -- Senioren 65+}}\\[2cm]
    {\Huge\textcolor{fach}{\textbf{Langentwurf}}}\\[0.5cm]
    {\large Tier 2 | Profil S}\\[2cm]
    {\LARGE\textbf{002 -- Meine Geräte-Inventur}}\\[0.5cm]
    {\large Modul A: Geräte \& Zugänge}\\[2cm]
    \begin{tabular}{rl}
        \textbf{Zeitrahmen:} & 45--60 Minuten \\
        \textbf{Vorkenntnisse:} & Einheit 001 (iPhone-Grundlagen) \\
    \end{tabular}
    \vfill
    {\normalsize Stand: \today}
\end{titlepage}

\tableofcontents
\newpage

% ============================================================================
\section{Bedingungsanalyse}

\subsection{Ausgangssituation}
\begin{itemize}
    \item Teilnehmer haben oft mehrere Apple-Geräte (iPhone, iPad, evtl. Mac)
    \item Wissen häufig nicht, welches Modell sie besitzen
    \item Können bei Support-Anfragen keine Geräteinformationen nennen
    \item Seriennummer ist unbekannt -- problematisch bei Verlust/Diebstahl
\end{itemize}

\subsection{Typische Probleme}
\begin{itemize}
    \item ``Ich weiß nicht, welches iPhone ich habe''
    \item ``Der Support fragt immer nach Dingen, die ich nicht weiß''
    \item ``Mein Enkel hat mir verschiedene Geräte eingerichtet''
    \item ``Sind auf allen Geräten die gleichen Konten?''
\end{itemize}

% ============================================================================
\section{Sachanalyse}

\subsection{Kernkonzepte}
\begin{enumerate}
    \item \textbf{Geräteidentifikation:} Modellname, Speicher, Seriennummer
    \item \textbf{Software-Version:} iOS, iPadOS, macOS -- wichtig für Kompatibilität
    \item \textbf{Apple-ID-Zuordnung:} Welches Konto auf welchem Gerät?
    \item \textbf{Dokumentation:} Sichere Aufbewahrung der Informationen
\end{enumerate}

\subsection{Praktische Relevanz}
\begin{tabular}{|l|p{8cm}|}
\hline
\textbf{Situation} & \textbf{Benötigte Information} \\
\hline
Support-Anruf & Modell, iOS-Version, Seriennummer \\
Verlust/Diebstahl & Seriennummer für Polizei \\
Verkauf & Modell, Speicher, Zustand \\
App-Kompatibilität & iOS-Version \\
iCloud-Probleme & Welche Apple-ID ist angemeldet? \\
\hline
\end{tabular}

% ============================================================================
\section{Didaktische Analyse}

\subsection{Stellung in der Reihe}
Einheit 002 folgt direkt auf 001 (iPhone-Grundlagen). Die Teilnehmer kennen jetzt die drei Codes und sollen nun ihre Geräte systematisch dokumentieren.

\subsection{Legitimation}
\begin{itemize}
    \item \textbf{Praktisch:} Jeder Support-Kontakt erfordert Geräteinformationen
    \item \textbf{Sicherheit:} Seriennummer bei Verlust wichtig
    \item \textbf{Übersicht:} Wissen, was man besitzt und wo welche Apple-ID aktiv ist
\end{itemize}

\subsection{Didaktische Reduktion}
\textbf{Ausgeklammert:} IMEI, Modellnummer (vs. Modellname), technische Spezifikationen (Chip, RAM).

\textbf{Fokus:} Nur die Informationen, die praktisch relevant sind.

% ============================================================================
\section{Lernziele}

\subsection{Grobziel}
Die Teilnehmer dokumentieren ihre Apple-Geräte vollständig und wissen, wo sie diese Informationen finden.

\subsection{Feinziele}
\begin{tabular}{|c|p{7cm}|c|c|}
\hline
\textbf{Nr.} & \textbf{Die Teilnehmer können...} & \textbf{Stufe} & \textbf{Niveau} \\
\hline
FZ1 & den Modellnamen ihres iPhones nennen & Wissen & Basis \\
FZ2 & die iOS-Version ihres Geräts finden & Anwenden & Basis \\
FZ3 & die Seriennummer lokalisieren & Anwenden & Basis \\
FZ4 & prüfen, welche Apple-ID angemeldet ist & Anwenden & Basis \\
FZ5 & erklären, warum diese Informationen wichtig sind & Verstehen & Vertiefung \\
\hline
\end{tabular}

% ============================================================================
\section{Methodische Analyse}

\subsection{Vorgehen}
\begin{enumerate}
    \item \textbf{Motivation:} Support-Szenario durchspielen (``Welches Gerät haben Sie?'')
    \item \textbf{Demonstration:} Am eigenen Gerät zeigen, wo die Infos stehen
    \item \textbf{Selbstständig:} Formular ausfüllen mit Anleitung
    \item \textbf{Kontrolle:} Gemeinsam prüfen, ob alles vollständig ist
\end{enumerate}

\subsection{Material}
Das vorhandene Material (Geraete-Inventur.pdf) enthält:
\begin{itemize}
    \item Seite 1: Warum \& Wie (Motivation + Anleitung)
    \item Seite 2: Inventar-Formulare (iPhone, iPad, Mac, Watch)
    \item Seite 3: Checkliste \& Tipps
\end{itemize}

% ============================================================================
\section{Verlaufsplanung}

\begin{longtable}{|p{1cm}|p{2cm}|p{5cm}|p{3cm}|p{2cm}|}
\hline
\textbf{Zeit} & \textbf{Phase} & \textbf{Inhalt} & \textbf{TN-Aktivität} & \textbf{Material} \\
\hline
0--5 & Einstieg & Support-Szenario: ``Welches iPhone haben Sie?'' & Nachdenken & -- \\
\hline
5--15 & Erklärung & Wo finde ich Modell, Version, Seriennummer? Am Gerät zeigen. & Mitmachen & S.1 \\
\hline
15--35 & Übung & Inventar-Formulare ausfüllen (eigene Geräte) & Selbst durchführen & S.2 \\
\hline
35--45 & Sicherung & Checkliste durchgehen, Apple-ID prüfen & Abhaken & S.3 \\
\hline
\end{longtable}

% ============================================================================
\section{Lösungen}

\textbf{Wo finde ich die Informationen?}
\begin{itemize}
    \item \textbf{iPhone/iPad:} Einstellungen $\rightarrow$ Allgemein $\rightarrow$ Info
    \item \textbf{Apple-ID:} Einstellungen $\rightarrow$ [Dein Name] oben
    \item \textbf{Mac:} Apfel-Menü $\rightarrow$ Über diesen Mac
\end{itemize}

% ============================================================================
\section{Didaktik-Score}

\begin{scorebox}
\begin{tabular}{|c|l|c|c|p{3.5cm}|}
\hline
\# & Dimension & Gew. & Score & Notiz \\
\hline
1 & Differenzierung & 15\% & 7/10 & Nur iPhone-Fokus, andere Geräte optional \\
2 & Taxonomie & 10\% & 9/10 & Anwenden dominiert (passend) \\
3 & Alltagsrelevanz & 10\% & 10/10 & Direkt praktisch nutzbar \\
4 & Aktivierung & 10\% & 10/10 & TN füllen eigene Daten aus \\
5 & Scaffolding & 10\% & 9/10 & Schritt-für-Schritt-Anleitung \\
6 & Feedback & 10\% & 8/10 & Checkliste vorhanden \\
7 & Sprachliche Klarheit & 10\% & 9/10 & Einfach, klare Pfade \\
8 & Motivation & 10\% & 9/10 & Support-Szenario motiviert \\
9 & Barrierefreiheit & 10\% & 8/10 & Formulare könnten größer sein \\
10 & Praktikabilität & 5\% & 9/10 & 45 Min realistisch \\
\hline
\multicolumn{3}{|r|}{\textbf{GESAMT:}} & \textbf{8.85/10} & \textcolor{warningorange}{Iteration empfohlen} \\
\hline
\end{tabular}
\end{scorebox}

\subsection{Verbesserungsvorschläge}
\begin{enumerate}
    \item \textbf{Differenzierung:} Separate Formulare für Basis (nur iPhone) vs. Vertiefung (alle Geräte)
    \item \textbf{Barrierefreiheit:} Formularfelder vergrößern, mehr Platz zum Schreiben
\end{enumerate}

\end{document}
